\section{Data and empirical design}\label{sec:empirical-design}

The preceding theoretical analysis in \Cref{sec:methodology} derives
restrictions that link distances between firms' latent characteristic
distributions to systematic covariance under maintained transmission
assumptions.
Those characteristic distributions and the transmission map are not observed.
The empirical goal is therefore modest: rather than recover the theoretical
objects, we ask whether an article-embedding proxy for the characteristic
geometry is associated with return co-movement in the theory-motivated
direction under the maintained empirical bridge.
The design proceeds in three steps.
First, for each firm, we construct an empirical distribution of embeddings
from selected financial news articles.
Second, we compute pairwise quadratic Wasserstein distances between
those distributions.
Third, we relate article-embedding distance to return co-movement in
a dyadic design.
Throughout, embeddings measure selected semantic content rather than latent
economic characteristics directly.
\Cref{sec:data-provenance} gives the full retrieval, filtering, and
provenance protocol.

\subsection{Hypothesis and testable
restrictions}\label{sec:testable-restrictions}

The theory delivers a conditional restriction on feasible covariance rather
than a point prediction, so the first exercise asks whether observed return
geometry is compatible with the envelope at declared transmission inputs.
The covariance envelope in \Cref{sec:methodology} brackets latent systematic
covariance through the characteristic-side distance and the maintained inputs
$(L,\tau)$.
To express that restriction in observed returns, we maintain centered factor
innovations, zero residual means, zero factor--residual cross covariances,
and zero cross-firm residual covariance under the same coupling $\pi$ used
for the systematic exposures.
With standardized total returns and unit marginal variances, these
restrictions identify $\rho_{ij}=\kappa_\pi$ and hence
$c^{\mathrm{actual}}_{ij}:=\E(R_i-R_j)^2=2-2\rho_{ij}$.
Writing $\Delta_{ij}$ for the transport excess under the realized coupling,
the observable classification algebra is
\begin{equation}\label{eq:random-transport-excess-bracket}
  c^{\mathrm{actual}}_{ij}-(w^+_{ij})^2
  \le \Delta_{ij}
  \le c^{\mathrm{actual}}_{ij}-(w^-_{ij})^2.
\end{equation}
A dyad is \emph{covered} when the lower endpoint is nonnegative,
\emph{violated} when the upper endpoint is negative, and
\emph{unresolved} otherwise.
The $(L,\tau)$ values are declared scenario inputs, so the classification
does not form a return-based estimate of either transmission input.
The observed article distance is only a proxy for an unobserved
characteristic distance.
Neither this inequality classification nor the dyadic regression below
identifies the transmission model.

The classification exercise evaluates compatibility at specified transmission
inputs, whereas the directional expectation can also be examined without
assigning each dyad to an envelope category.
The second empirical check is therefore a reduced-form regression of return
distance on article-embedding distance after absorbing stable
firm-level heterogeneity.
Its outcome is the return chord distance
\begin{equation}\label{eq:chord-distance}
  \hat{D}_{ij} = \sqrt{2(1 - \hat{\rho}_{ij})},
\end{equation}
where $\hat{\rho}_{ij}$ is the sample Pearson correlation of daily returns
(equation~\ref{eq:sample-correlation} below).
On this scale, $D_{ij}=0$ denotes perfectly correlated returns,
$D_{ij}=\sqrt{2}$ denotes zero correlation, and $D_{ij}=2$ denotes perfectly
negative correlation.
Larger chord distance therefore means weaker return co-movement.

The regression is defined over the observed undirected pairs $i<j$:
\begin{equation}\label{eq:dyadic-fixed-effects}
  D_{ij}=\mu+a_i+a_j+\beta_X X_{ij}+z_{ij}^{\top}\gamma+u_{ij},
  \qquad 1\le i<j\le N.
\end{equation}
Here $N=\ppvalue{n-firms}$, $\mu$ is the intercept, $X_{ij}$ is the
article-embedding regressor with coefficient $\beta_X$, $z_{ij}$ is the
vector of secondary controls with coefficient vector $\gamma$, and $u_{ij}$
is the dyadic error.
The terms $a_i+a_j$ are symmetric additive firm effects.
The restriction $i<j$ counts each unordered pair once.
The normalization omits one firm-effect indicator; changing the omitted
ticker does not change $\widehat\beta_X$.
The canonical specification sets $X_{ij}=W^{(2)}_{ij}$ and $z_{ij}=0$,
yielding the fixed-W$_2$ slope $\beta_{W_2}$.

\begin{hypothesis}[Reduced-form distance association]
  \label{hyp:monotonicity}
  For the theory-aligned arm, the reported two-sided regression test is
  \begin{equation}\label{eq:dyadic-association-hypothesis}
    H_0:\beta_{W_2}=0
    \qquad\text{against}\qquad
    H_1:\beta_{W_2}\ne0.
  \end{equation}
  The directional expectation, stated before inspecting the result, is
  $\beta_{W_2}>0$: under the maintained bridge, firms farther apart in the
  characteristic geometry should tend to be farther apart in
  return-correlation space after symmetric firm effects.
  This directional expectation motivates the reported two-sided test but is
  not a causal claim.
  The envelope diagnostic separately asks whether the nonnegative
  transport-excess restriction is covered, violated, or unresolved at each
  declared $(L,\tau)$ scenario; it is a conditional inequality
  classification, not a fitted-parameter test.
\end{hypothesis}

\subsection{Sample and return construction}\label{sec:sample-and-window}

Firm-level distribution comparisons require sufficient text and common
coverage across the fixed window to construct a common-size empirical measure
for every firm, which motivates a coverage-conditioned sample rather than a
broad but uneven news universe.
The news universe is a prespecified, Nasdaq-100-based frame of
\ppvalue{n-firms-news} Nasdaq-listed firms.
Each retained firm had at least \ppvalue{annual-raw-article-floor} raw articles
in every calendar year from
\ppvalue{sample-year-start} through \ppvalue{sample-year-end} in Nasdaq's
public per-symbol news archive.
Within the fixed frame, articles satisfying the date-window and
minimum-body-length filters are ordered by URL hash, and up to
\ppvalue{n-samples-cap} are retained over the pooled window.
The common coverage rule improves comparability across firm-level empirical
measures, but the resulting panel is not representative of all
Nasdaq-listed firms.
\Cref{tab:sample-descriptives-sectors} reports the sector composition,
per-firm article counts, return volatilities, and within- versus cross-sector
W$_2$ distances.
\Cref{tab:sample-flow} records the resulting flow from source records to
estimation dyads; the full retrieval and provenance protocol is in
\Cref{sec:data-provenance}.

% AUTO-GENERATED by pipeline render-paper1-tables -- do not edit by hand.

\begin{table}[H]

\centering

\small
\begin{tabularx}{\linewidth}{Xrrrrr}
\toprule
Sector & $n$ & Articles & Volatility & $W_2$ within & $W_2$ cross \\
\midrule
Communication Services & 11 & 128.0 & 0.0239 & 1.042 & 1.060 \\
Consumer Cyclical & 18 & 128.0 & 0.0281 & 1.033 & 1.053 \\
Consumer Defensive & 7 & 128.0 & 0.0176 & 1.010 & 1.058 \\
Energy & 1 & 128.0 & 0.0386 & \textemdash & 1.045 \\
Financial Services & 1 & 128.0 & 0.0272 & \textemdash & 1.070 \\
Healthcare & 14 & 128.0 & 0.0233 & 1.025 & 1.065 \\
Industrials & 7 & 128.0 & 0.0240 & 1.030 & 1.057 \\
Technology & 38 & 128.0 & 0.0265 & 0.998 & 1.050 \\
Utilities & 3 & 128.0 & 0.0159 & 0.882 & 1.058 \\
\midrule
All sectors & 100 & 128.0 & 0.0251 & 1.008 & 1.056 \\
\bottomrule
\end{tabularx}

\caption{Per-sector sample composition: firm count, mean article count per firm, mean daily return volatility, and mean within- versus cross-sector pairwise rooted $W_2$ distance. Articles are subsampled to a fixed per-firm cap, so counts are near-constant by construction. A dash indicates a statistic that is undefined for that sector (e.g.\ within-sector distance for a singleton sector). Authors' calculations.}

\label{tab:sample-descriptives-sectors}

\end{table}


Two features of that composition motivate later design choices.
First, mean daily return volatility differs markedly across sectors, whereas
the primary return chord distance is constructed from correlations and is
therefore invariant to firm-specific return scale.
Symmetric firm effects instead absorb stable endpoint heterogeneity.
Second, mean within-sector W$_2$ distance is smaller than the corresponding
cross-sector distance in every sector holding more than one firm, so
article-embedding distance encodes sector information before any return is used.
The sector diagnostic in \Cref{sec:geometric-representation} quantifies that
pattern, and the pooled same-sector rung probes it.
The canonical symmetric-firm-effects specification does not include a
same-sector control; the pooled rung shows how much of the unadjusted
association survives that descriptive comparison.

To measure co-movement independently of firm-specific return scale, the
outcome is constructed from correlations of daily simple returns.
Daily simple returns are Yahoo Finance adjusted-close ratios minus one.
Return-dependent tests use the exact complete-case matrix of
\ppvalue{return-panel-n-observations} dates shared by
\ppvalue{return-panel-n-tickers} priced firms, spanning
\ppvalue{return-panel-date-start} through \ppvalue{return-panel-date-end}.
Nullable observations are removed before matrix construction, with no
interpolation or zero filling.
The news and return panels overlap over the
\ppvalue{sample-year-start}--\ppvalue{sample-year-end} calendar window, so
the association is contemporaneous and reduced-form rather than a
pre-period prediction.
Return alignment and sector-metadata details are documented in
\Cref{sec:data-provenance}.

The sample correlation underlying the chord distance in
equation~\ref{eq:chord-distance} is
\begin{equation}\label{eq:sample-correlation}
  \hat{\rho}_{ij}
  =
  \frac{
    \sum_{t=1}^{N_T} (r_t^i - \bar{r}^i)(r_t^j - \bar{r}^j)
  }{
    \sqrt{
      \sum_{t=1}^{N_T} (r_t^i - \bar{r}^i)^2
      \sum_{t=1}^{N_T} (r_t^j - \bar{r}^j)^2
    }
  },
\end{equation}
where $\bar{r}^i$ is the sample average of asset $i$'s returns, yielding a
$\ppvalue{n-firms} \times \ppvalue{n-firms}$ symmetric chord distance matrix
$D$ with zero diagonal.

\subsection{Article-embedding representation and distance}
\label{sec:distance-construction}\label{sec:article-proxy-model}

The empirical bridge requires a common representation of heterogeneous
articles while retaining each firm as a distribution rather than collapsing
its text to a single point.
Embeddings provide that measurement technology by mapping documents into a
shared numerical space in which semantic proximity can be compared.
An embedding model is a neural network that maps a text document to a
fixed-length numerical vector---its \emph{embedding}.
Documents discussing similar topics are placed nearby in the resulting
high-dimensional space, whereas documents on unrelated subjects are placed
farther apart.
A firm's collection of vectors can therefore be treated as an empirical
distribution in semantic space.
This representation measures selected semantic content; it does not observe
the firm's latent economic characteristics directly.

Wasserstein comparison further requires each article to occupy the same
reusable geometry before firm pairs are formed.
A \emph{bi-encoder} meets that requirement by processing each document
independently into one reusable vector, so every within- and cross-firm
comparison shares a common geometry.
A cross-encoder instead returns pair-specific scores and would not directly
define the per-firm empirical measures needed for Wasserstein computation.
The application therefore uses the bi-encoder as a fixed text map, without
retrieval or reranking.

The distinction between the theoretical object and the article-embedding
proxy can be stated formally.
The randomized framework in \Cref{sec:methodology} operates on latent
characteristic laws $C_i\in\mathcal P_2(\Omega)$ over an abstract risk-factor
space $\Omega$.
Market data do not reveal these laws.
In the application, write $C_i^\ast$ for the latent economic characteristic
distribution represented by the theoretical $C_i$, and let $Q_i$ denote the
distribution of content selected for media coverage and mapped by the
embedding function $\phi$ into semantic space.
The observed articles define only the empirical embedding measure
\begin{equation}\label{eq:empirical-embedding-measure}
  \widehat Q_i
  :=\frac{1}{N_i}\sum_{d=1}^{N_i}\delta_{\phi(d)}.
\end{equation}
The data therefore measure distances between $\widehat Q_i$, not between the
latent $C_i^\ast$.
Reading the geometry of $Q_i$ as a proxy for that of $C_i^\ast$ is the
maintained measurement premise below, not an identity between semantic and
economic characteristics.

\begin{assumption}[Measurement Premise: Article-Embedding
  Proxy]\label{ass:article-proxy}
  After symmetric firm margins and the declared controls, distortions from
  media selection, article generation, duplicate handling, text extraction,
  and embedding do not covary with return distance strongly enough to create
  or reverse the sign of the latent association.
\end{assumption}

For the randomized-exposure formulation, the application also maintains that
the same bi-Lipschitz map and common $L$ apply across firms and that the
latent characteristic laws have finite second moments.
Under \Cref{ass:article-proxy}, the observed embedding law $\widehat Q_i$
proxies the latent characteristic law $C_i^\ast$.
The common randomized map then sends $C_i^\ast$ to the latent exposure law
$P_i$, while $\tau_i$ bounds only transmission deviation in risk coordinates,
not article-measurement error.
These are scenario assumptions, not fitted measurement equations.
Article and return data identify neither $P_i$, $L$, nor $\tau_i$.

The proxy is motivated by financial news because articles repeatedly record
analysts', journalists', and market participants' assessments of
firm-specific risks.
Across the window, article topics cover persistent firm attributes, such as
industry structure and business model, and transient events, such as earnings
surprises and management changes.
Accordingly, $C_i^\ast$ and $Q_i$ are pooled distributions over the fixed
sample window, not timeless structural characteristics, and the corresponding
return object is the same-window covariance.
Treating the exposure law $P_i$ as fixed over the window is therefore a
pooled-window approximation.
The design does not separately identify the media-and-embedding distortion in
\Cref{ass:article-proxy}.

Given this proxy role, the headline representation embeds full article bodies
with the open-weights \texttt{qwen/qwen3-embedding-8b} bi-encoder, producing
\ppvalue{embed-dim}-dimensional vectors under pinned serving parameters
\citep{qwen3_embedding_2025}.
Each news article $d$ associated with asset $i$ is mapped to an embedding
vector $\bm{e}_d=\phi(d)\in\R^{\ppvalue{embed-dim}}$, and these vectors are
the atoms of $\widehat Q_i$ in equation~\ref{eq:empirical-embedding-measure}.
Their distances approximate differences in selected semantic content, subject
to \Cref{ass:article-proxy}.
Every article predates the later returns evaluated in
\Cref{sec:post-window-dyadic}, but the encoder has no documented
training cutoff.
Because the encoder is trained for semantic retrieval rather than recovery of
financial risk factors, its economic interpretation rests on the measurement
premise in \Cref{ass:article-proxy}.

Holding article counts fixed keeps the number and weight of atoms comparable
across firms rather than combining semantic comparisons with unequal
empirical support sizes.
Within each firm, articles are ordered lexicographically by their
stable URL hash.
All full-width encoder comparisons use the first $m=128$ articles under this
ordering; for the headline encoder we additionally use $m\in\{32,64,128\}$.
Let $x_{ia}$ denote the row-normalized embedding of selected article $a$, and
let $S_m$ be the set of permutations of $\{1,\ldots,m\}$.
The primary statistic is
\begin{equation}\label{eq:empirical-wasserstein-2}
  \widehat W_{2,m}(i,j)
  = \left(\min_{\pi\in S_m}\frac{1}{m}\sum_{a=1}^{m}
  \norm{x_{ia}-x_{j,\pi(a)}}_2^{2}\right)^{1/2}.
\end{equation}
The minimization pairs every article from one firm with one article from the
other and finds the smallest root-mean-square semantic displacement, so W$_2$
is small when two firms' article topics can be aligned closely even if no
individual article is shared.
The fixed-W$_2$ diagnostic uses the common $m=128$ unit-chord Qwen3
representation.
The complete W$_1$ and energy definitions are in
Appendix~\ref{sec:representation-diagnostics} as lower-cost descriptive
comparators, and the theorem is agnostic to the empirical ground metric.

Before relating the article-embedding geometry to returns, we ask a
descriptive question: does the pairwise W$_2$ matrix exhibit economically
recognizable organization?
The visualization is intuition-building scaffolding, not a test of the theory
or the measurement premise.
Because pairwise distances are not themselves points in a vector space,
metric multidimensional scaling (MDS) recovers an approximate map from the
distance table alone---analogous to reconstructing a city layout from a table
of driving distances.
\Cref{fig:sector-mds} projects the exact W$_2$ matrix into two dimensions and
colours each firm by its Nasdaq summary sector label.

\begin{figure}[htbp]
  \centering
  \ppfigure{0.96}{mds_sector_plot}
  \caption{Descriptive two-dimensional metric MDS projection of the exact
    128-article W$_2$ matrix.
    Both axes are unitless projection coordinates; each labelled point is a
    priced firm, and colours follow the sector legend.
    Sector colours overlap throughout the map rather than separating into
    discrete groups.
    All statistics and pair selections use the original W$_2$ distances, not
    these approximate coordinates, whose local-neighbour fidelity is limited
    (\Cref{fig:mds-shepard-diagnostic}).
  Authors' calculations.}
  \label{fig:sector-mds}
\end{figure}

In finance terms, the overlap suggests that the article-embedding geometry
contains broad sector structure but is not exhausted by sector labels.
The DISCO (distance components) decomposition quantifies this distinction:
sector membership accounts for only $R^2_{\mathcal E}=\ppnum[3]{disco-r2}$ of
total energy dispersion \citep{szekely_energy_2013}.
DISCO partitions dispersion in the distance matrix into between- and
within-sector shares, much as a one-way ANOVA partitions variance, but it
operates directly on pairwise distances rather than on a scalar variable.
The full DISCO decomposition and Shepard-fidelity checks are reported in
\Cref{sec:geometric-representation}.
This descriptive structure neither validates the theory nor establishes that
embeddings recover latent characteristics.
It instead motivates the question in \Cref{sec:post-window-dyadic}: whether
firm-level article-embedding differences in the original W$_2$ matrix are
associated with return distance after symmetric firm effects.
Having defined the sample, return outcome, and article-embedding geometry, we
now specify the estimator and inference used to answer that question.

\subsection{Estimator and inference}
\label{sec:primary-estimator}\label{sec:empirical-testing}
\label{sec:inference-design}

Article and return data do not identify $L$ or $\tau$, so the envelope check
treats both as sensitivity inputs rather than fitted parameters.
The primary random-exposure check evaluates the one-sided transport-excess
interval in
equations~\ref{eq:random-w2-lower-metric}--\ref{eq:random-kappa-bracket}.
It treats the fixed Qwen3 $m=128$ W$_2$ matrix defined in
\Cref{sec:distance-construction} as the characteristic-side distance $G_{ij}$
and uses standardized total-return geometry
$c^{\mathrm{actual}}_{ij}=2-2\hat\rho_{ij}$.
The nine $(L,\tau)$ cells form an uncalibrated local sensitivity grid around
isometry and zero slack: $L\in\{1,1.25,1.5\}$ and
$\tau_i=\tau_j\in\{0,0.05,0.10\}$.

Dyads sharing a firm may be dependent, so the cross-sectional information is
governed by the number of firms rather than the number of pairs.
Firms are therefore the resampling units, and symmetric additive firm effects
absorb stable endpoint heterogeneity.
Each multinomial node-bootstrap draw assigns a count to every firm and
weights dyad $(i,j)$ by the product of its endpoint counts, handling
dependence among dyads that share a firm but not pair-specific
omitted variables.
Percentile intervals and the finite-draw-corrected two-sided sign tail use
the same 1,999-draw schedule for every reported specification.

The residual representation, partial-$R^2$ definition, nuisance-rank details,
six-rung estimator audit, exact weights, rank checks, functional-form
diagnostics, and influence results are reported in
Appendix~\ref{sec:dyadic-diagnostics}.

The contemporaneous design cannot distinguish the association from conditions
shared by the article and return windows, so a timing-separated specification
asks whether the qualitative relation persists when the return window
moves forward.
It holds the \ppvalue{sample-year-start}--\ppvalue{sample-year-end}
article-embedding W$_2$ matrix fixed and rebuilds the outcome and symmetric
firm effects from the 2023--2026 complete-case return panel.
It uses the same dyads, estimator, bootstrap schedule, and two-sided null.
This separation prevents returns in the evaluation window from determining
article selection.
It is nevertheless timing-separated evidence rather than point-in-time
prediction: the encoder is not a point-in-time model, and common latent
drivers may persist across windows.

With the primary design fixed, the final subsection describes how the
analysis probes sensitivity to the chosen article-embedding representation.

\subsection{Representation sensitivity ladder}\label{sec:representation-ladder}

Because embeddings are a measurement choice, the return-distance association
would be less informative if it appeared only under one model, scale, output
width, metric, or training-corpus vintage.
The representation sensitivity ladder therefore asks whether the estimate
survives when each of these axes is varied while the empirical design
remains fixed.
It holds the estimator, return matrix, firm sample, and node-bootstrap
schedule fixed while changing encoder, output width, or distance metric,
thereby probing sensitivity to these design choices
(\Cref{tab:w2-representation-sensitivity}).

The seven-row base ladder targets three distinct measurement concerns.
The capacity comparison asks whether model scale drives the association by
pairing Qwen3-Embedding-8B with its 4B variant at the same native width and
at a common 1,024-dimensional output.
The output-width comparison asks whether the estimated association depends on
the full coordinate set.
It uses Matryoshka prefixes, for which the encoder's training makes
leading-coordinate prefixes usable as nested, lower-capacity embeddings at
256 and 64 coordinates.
These prefixes are fixed before the analysis, not principal components or
selected coordinates \citep{kusupati_matryoshka_2022}.
The architecture comparison asks whether the association is specific to the
Qwen training pipeline.
It uses BGE-large-v1.5, a bidirectional encoder from a different provider and
training pipeline that changes architecture, pretraining, and provider
jointly \citep{xiao_cpack_2023}.
The compact contrast table reports the candidate-minus-reference pairs
descriptively with paired intervals and Holm-adjusted $p$-values; none is a
causal or encoder-superiority claim.
W$_1$ and energy are substituted for $X_{ij}$ only in the descriptive
representation table.
Because all rows share the node schedule, endpoint effects, and omitted
indicator, their standardized effects are comparable even though one-unit
coefficients inherit different distance scales.

Training-data vintage requires a separate design because comparisons across
off-the-shelf encoders also change architecture and training pipeline.
The matched EttaX panel holds those features fixed so that only the Wikipedia
snapshot varies.
Its three arms are 12-layer, 320-dimensional Transformer bi-encoders with
approximately 22~million parameters, trained by the authors on Wikipedia
article text using a masked autoencoding objective.
During training, 60\% of input tokens are randomly deleted and a separate
single-layer decoder reconstructs the original document from the encoder's
compressed representation.
Architecture, optimizer, training budget (900~million token slots),
byte-level BPE tokenizer, corpus sampler, and random seed are identical
across arms; only the Wikipedia snapshot changes.
V0 uses the 20 December 2017 snapshot, which predates the
\ppvalue{sample-year-start}--\ppvalue{sample-year-end} article sample window.
V1 uses the 20 December 2020 snapshot, which falls within the sample window.
V3 uses an August 2026 snapshot, which postdates the return evaluation window.
If the measured association were an artifact of temporal overlap between
encoder training data and the article sample, V0 would show a weaker
association than V1 or V3.
The vintage equivalence test declares equivalence only when the full paired
interval lies inside a symmetric margin calibrated to the previously computed
Qwen3-8B--BGE-large contrast.
These margins are benchmark-calibrated robustness thresholds, not causal or
architecture-only comparisons.
These exercises assess the stability of the article-embedding proxy
association; they do not identify the transmission model or convert the
design into point-in-time prediction.

% Geometric diagnostics (MDS, DISCO) are woven into §4.4 at point of use.
% Labels preserved for any inbound cross-references.
\label{sec:geometric-diagnostics}
